\chapter{গবেষণা পদ্ধতি}\label{ch:method}

\section{গবেষণার দৃষ্টিভঙ্গি}
এই গবেষণা \emph{ডিজাইন-সায়েন্স} ও \emph{পরিমাণগত} দৃষ্টিভঙ্গি অনুসরণ করে। এর প্রধান ফলাফল একটি নির্মিত বস্তু, অর্থাৎ অ্যানালিটিক্স
কাঠামো, যা একটি নির্দিষ্ট ব্যবসায়িক সমস্যা সমাধানের জন্য তৈরি এবং পরে সুস্পষ্ট মানদণ্ডে কঠোরভাবে মূল্যায়িত। অভিজ্ঞতামূলক বিশ্লেষণ
পরিমাণগত এবং বিক্রয় লেনদেনের একটি বড় মাধ্যমিক ডেটাসেট ব্যবহার করে। সামগ্রিক জীবনচক্র CRISP-DM \citep{shearer2000} অনুসরণ
করে, যা প্রকল্প প্রস্তাবনায় দশটি ধাপে পরিমার্জিত হয়েছে (\cref{fig:lifecycle})।

\begin{figure}[htbp]\centering
\begin{tikzpicture}[node distance=4mm and 4mm, every node/.style={font=\small},
  box/.style={draw=brand, rounded corners, fill=brand!8, minimum width=3.3cm, minimum height=9mm, align=center},
  arr/.style={-{Stealth[length=2mm]}, thick, brand}]
\node[box] (p1) {১. সমস্যা\\সংজ্ঞায়ন};
\node[box, right=of p1] (p2) {২. ডেটা\\সংগ্রহ};
\node[box, right=of p2] (p3) {৩. ডেটা\\প্রস্তুতি};
\node[box, right=of p3] (p4) {৪. অনুসন্ধানমূলক\\বিশ্লেষণ};
\node[box, below=9mm of p4] (p5) {৫. ফিচার\\ইঞ্জিনিয়ারিং};
\node[box, left=of p5] (p6) {৬. মডেল\\উন্নয়ন};
\node[box, left=of p6] (p7) {৭. মডেল\\মূল্যায়ন};
\node[box, left=of p7] (p8) {৮. ড্যাশবোর্ড\\উন্নয়ন};
\node[box, below=9mm of p8] (p9) {৯. সিস্টেম\\মূল্যায়ন};
\node[box, right=of p9] (p10) {১০. ডকুমেন্টেশন};
\draw[arr] (p1)--(p2); \draw[arr] (p2)--(p3); \draw[arr] (p3)--(p4); \draw[arr] (p4)--(p5);
\draw[arr] (p5)--(p6); \draw[arr] (p6)--(p7); \draw[arr] (p7)--(p8); \draw[arr] (p8)--(p9); \draw[arr] (p9)--(p10);
\draw[arr, dashed, accent] (p7.north) to[bend left=20] node[above, font=\scriptsize]{পুনরাবৃত্তি} (p3.south);
\end{tikzpicture}
\caption{প্রকল্পের জীবনচক্র, CRISP-DM থেকে অভিযোজিত}\label{fig:lifecycle}
\end{figure}

\section{সমস্যা সংজ্ঞায়ন}
ব্যবসায়িক প্রশ্নগুলো থেকে চারটি বিশ্লেষণী সমস্যা নির্ধারণ করা হয়েছে:
\begin{enumerate}
  \item \textbf{চাহিদা পূর্বাভাস (রিগ্রেশন)।} পরের মাসে প্রতিটি পণ্যের কত ইউনিট বিক্রি হবে তা পূর্বাভাস করা। এটি উৎপাদন ও মজুদ
        পরিকল্পনায় সহায়তা করে।
  \item \textbf{বড় অর্ডারের পূর্বাভাস (ক্লাসিফিকেশন)।} অর্ডার এন্ট্রির মুহূর্তে অনুমান করা অর্ডারটি বড় হবে কি না (\(\geq 10\) ইউনিট)।
        এটি মজুদ ও ডেলিভারি সক্ষমতা বণ্টনে সহায়তা করে।
  \item \textbf{বিভাজন (ক্লাস্টারিং)।} একই রকম আচরণের পণ্য ও বাজার-কোষ (অঞ্চল $\times$ চ্যানেল $\times$ গ্রাহকের ধরন) দলবদ্ধ করা,
        যা পোর্টফোলিও ও বিক্রয়-বাহিনী কৌশলে সহায়তা করে।
  \item \textbf{অ্যানোমালি ডিটেকশন।} অস্বাভাবিক ব্যবসায়িক দিন ও সন্দেহজনক লেনদেন শনাক্ত করা, যা পর্যবেক্ষণ ও ডেটার মান নিয়ন্ত্রণে
        সহায়তা করে।
\end{enumerate}

\section{ডেটা}
ডেটাসেটে জানুয়ারি 2020 থেকে সেপ্টেম্বর 2026 পর্যন্ত \NCols টি বৈশিষ্ট্যসহ \NRows টি লেনদেন রয়েছে। প্রতিটি রেকর্ডে আছে একটি শনাক্তকারী ও
তারিখ; পণ্য (শনাক্তকারী, নাম, ক্যাটাগরি, প্যাক সাইজ, চালুর বছর); মূল্য-তথ্য (মুদ্রিত সর্বোচ্চ খুচরা মূল্য, তালিকা মূল্য, বাজার মূল্য,
ডিসকাউন্টের শতাংশ ও একক বিক্রয়মূল্য); পরিমাণ; বাংলাদেশি টাকায় মোট, ডিসকাউন্ট ও নিট বিক্রয়মূল্য; এবং অঞ্চল, চ্যানেল ও গ্রাহকের
ধরন। পূর্ণ ডেটা অভিধান \cref{app:dictionary}-এর \cref{tab:dictionary}-এ দেওয়া আছে। ডেটায় কোনো ব্যক্তিগত শনাক্তকারী তথ্য নেই।

\section{মূল্যায়ন নকশা}
\subsection{সময়ভিত্তিক hold-out}
যেহেতু সমস্যাগুলো সময়ের সঙ্গে জড়িত, প্রতিটি supervised মডেল \textbf{2020--2024} সালে প্রশিক্ষিত এবং \textbf{জানুয়ারি
2025 -- সেপ্টেম্বর 2026}-এ পরীক্ষিত। \TestMonths মাসের এই অদেখা সময়কাল ``ভবিষ্যৎ''-এর ভূমিকা পালন করে। ইচ্ছাকৃতভাবে
এলোমেলো train/test বিভাজন ব্যবহার করা হয়নি, কারণ তাতে মডেল যে পর্যবেক্ষণগুলো পূর্বাভাস করছে তার পরের পর্যবেক্ষণ থেকেই শিখে
ফেলত \citep{bergmeir2012}।

\subsection{Hyper-parameter টিউনিং}
Hyper-parameter কেবল প্রশিক্ষণের বছরগুলোতে \emph{expanding-window} যাচাই দিয়ে টিউন করা হয়েছে। যাচাই-বছর
$y \in \{2022, 2023, 2024\}$-এর জন্য মডেলটি $y$-এর আগের সব মাসে প্রশিক্ষিত হয়ে $y$-তে স্কোর পায়। গড় যাচাই-স্কোরে সেরা
কনফিগারেশনটি এরপর 2020--2024-এ পুনরায় প্রশিক্ষিত হয়ে কেবল একবার পরীক্ষাকালে মূল্যায়িত হয়। ক্লাসিফিকেশনে probability
threshold-ও একইভাবে বেছে নেওয়া হয়েছে: মডেলটি 2020--2023-এ প্রশিক্ষিত হয়েছে এবং 2024-এ যে threshold F1 সর্বোচ্চ করে সেটি
রাখা হয়েছে।

\subsection{মেট্রিক}
রিগ্রেশনের জন্য নিচের পরিমাপগুলো ব্যবহৃত, যেখানে $n$টি পণ্য-মাসের জন্য $y_i$ প্রকৃত ও $\hat y_i$ পূর্বাভাসিত মান:
\begin{align*}
\text{MAE} &= \tfrac{1}{n}\textstyle\sum_i |y_i-\hat y_i|, &
\text{RMSE} &= \sqrt{\tfrac{1}{n}\textstyle\sum_i (y_i-\hat y_i)^2}, &
R^2 &= 1-\frac{\sum_i (y_i-\hat y_i)^2}{\sum_i (y_i-\bar y)^2},\\
\text{MAPE} &= \tfrac{100}{n}\textstyle\sum_i \frac{|y_i-\hat y_i|}{\max(y_i,1)}, &
\text{WAPE} &= 100\,\frac{\sum_i |y_i-\hat y_i|}{\sum_i y_i}. &&
\end{align*}
ব্যবসায়িক যোগাযোগে WAPE অগ্রাধিকার পায়, কারণ এটি প্রতিটি পণ্যকে তার পরিমাণ অনুযায়ী ভার দেয়। ক্লাসিফিকেশনে accuracy,
precision, recall, F1, ROC-AUC, PR-AUC ও শীর্ষ দশমাংশের lift উল্লেখ করা হয়েছে। ক্লাস্টারিংয়ে silhouette coefficient,
Davies--Bouldin সূচক ও Calinski--Harabasz সূচক। অ্যানোমালি ডিটেক্টরগুলো মূল্যায়িত হয়েছে জানা ধরনের রোপিত ত্রুটিসহ ডেটার একটি
অনুলিপিতে precision ও recall দিয়ে।

\subsection{বেঞ্চমার্ক}
\citet{makridakis2018} অনুসরণ করে প্রতিটি পূর্বাভাস মডেলকে দুটি naive নিয়মের সঙ্গে তুলনা করা হয়েছে: \emph{naive} (পরের মাস = এই
মাস) এবং \emph{seasonal naive} (পরের মাস = গত বছরের একই মাস)। কোম্পানি-পর্যায়ে Holt--Winters ও SARIMA পরিসংখ্যানিক বেঞ্চমার্ক।
ক্লাসিফিকেশনের বেঞ্চমার্ক সংখ্যাগরিষ্ঠ-শ্রেণি নিয়ম।

\section{পরিসংখ্যানিক বিশ্লেষণ}
অনুসন্ধানমূলক ধাপে ব্যবহৃত হয়েছে বর্ণনামূলক পরিসংখ্যান, Spearman rank correlation (প্রাপ্ত বিষম বণ্টনের প্রতি robust),
দলগুলোর পার্থক্য যাচাইয়ে Kruskal--Wallis টেস্ট, এবং স্বাধীনতা যাচাইয়ে chi-square টেস্ট। নমুনা অত্যন্ত বড় হওয়ায় প্রায় যেকোনো
পার্থক্যই পরিসংখ্যানিকভাবে তাৎপর্যপূর্ণ হয়। তাই \emph{effect size} ($\varepsilon^2$, $\rho^2$ ও Cram\'er's $V$) উল্লেখ করে
ব্যবহারিক গুরুত্ব বিচার করা হয়েছে। চাহিদার মূল্য-স্থিতিস্থাপকতা (price elasticity) নির্ণয়ে একটি fixed-effects log--log রিগ্রেশন
ব্যবহৃত হয়েছে।

\section{টুল ও পরিবেশ}
কাঠামোটি Python~3-এ বাস্তবায়িত: ডেটা পরিচালনায় pandas ও NumPy \citep{mckinney2022}, পরিসংখ্যান ও time-series মডেলে SciPy ও
statsmodels, মেশিন লার্নিংয়ে scikit-learn \citep{pedregosa2011}, gradient boosting-এ XGBoost \citep{chen2016}, ভিজ্যুয়ালাইজেশনে
Matplotlib ও Plotly, ডাইমেনশনাল ডেটাবেসে SQLite, এবং ওয়েব ড্যাশবোর্ডে Streamlit। প্রতিবেদনটি \LaTeX-এ লেখা। সব এলোমেলো
প্রক্রিয়ায় একটি নির্দিষ্ট seed (42) ব্যবহৃত, তাই প্রতিটি ফলাফল হুবহু পুনরুৎপাদন করা যায়।

\section{নৈতিক বিবেচনা}
ডেটাসেটে কোনো নাম, যোগাযোগের তথ্য বা অন্য কোনো ব্যক্তিগত শনাক্তকারী তথ্য নেই; গ্রাহকরা কেবল ধরন (যেমন \emph{Retailer}) ও
অঞ্চল দিয়ে বর্ণিত। প্রস্তাবনায় উল্লিখিত নৈতিক নীতি অনুযায়ী ডেটা কেবল শিক্ষাগত উদ্দেশ্যে ব্যবহৃত হয়েছে। মডেলের সীমাবদ্ধতা ও
অনিশ্চয়তা স্বচ্ছভাবে উল্লেখ করা হয়েছে (\cref{sec:limitations}), যাতে সিদ্ধান্ত গ্রহণে পূর্বাভাসের উপর অতিরিক্ত আস্থা না রাখা হয়।
